% Generated by docs/research_report/tables/make_tables.py. Do not edit by hand.
\begin{table}[htbp]
\centering
\caption[Baseline hyperparameters used in the reported runs]{Baseline hyperparameters from the saved run configurations. All use Adam, inverse-frequency class-weighted cross-entropy, a maximum of 300 epochs and validation patience 25. Width is the hidden dimension; decay is weight decay. Rare min. is the minimum category frequency for TE-G-SAGE, whose edge MLP width is 128. Published and refit E-GraphSAGE initialise nodes with ones; TE-G-SAGE uses a learned constant. The $+$ features ablations replace that initialisation with our ten node features. Published refers to the architectural settings, subject to deviations D1--D5 in the text.}
\label{tab:baseline-settings}
\small
\resizebox{\textwidth}{!}{%
\begin{tabular}{llrrrrrr}
\toprule
Dataset & Variant & Layers & Width & Dropout & Learning rate & Decay & Rare min. \\
\midrule
\multicolumn{8}{l}{\emph{E-GraphSAGE}} \\
UNSW & published & 2 & 128 & 0.2 & 0.001 & 0 & -- \\
UNSW & refit & 1 & 64 & 0.2 & 0.0005 & 0 & -- \\
UNSW & $+$ features & 1 & 64 & 0.2 & 0.0005 & 0 & -- \\
ToN & published & 2 & 128 & 0.2 & 0.001 & 0 & -- \\
ToN & refit & 2 & 128 & 0.2 & 0.001 & 0 & -- \\
ToN & $+$ features & 2 & 128 & 0.2 & 0.001 & 0 & -- \\
\midrule
\multicolumn{8}{l}{\emph{TE-G-SAGE}} \\
UNSW & published & 2 & 128 & 0.3 & 0.0003 & 0.0001 & 50 \\
UNSW & refit & 2 & 32 & 0.3 & 0.001 & 0.0001 & 2 \\
UNSW & $+$ features & 2 & 32 & 0.3 & 0.001 & 0.0001 & 2 \\
ToN & published & 2 & 128 & 0.3 & 0.0003 & 0.0001 & 50 \\
ToN & refit & 1 & 32 & 0.3 & 0.001 & 0.0001 & 2 \\
ToN & $+$ features & 1 & 32 & 0.3 & 0.001 & 0.0001 & 2 \\
\bottomrule
\end{tabular}
}
\end{table}
